\begin{statementbox}

	\textbf{\textcolor{CStatement}{条件}}
	\begin{description}[style=nextline, leftmargin=2.8em, labelsep=0.5em, itemsep=0.35em]
		\item[\textbf{\textcolor{CStatement}{条件 1（侧棱相等）：}}] 设 $A,B,C$ 不共线，$S\notin(ABC)$。三棱锥 $S$-$ABC$ 中，侧棱
		      \[
			      SA=SB=SC.
		      \]
		      等价地，顶点 $S$ 在底面 $ABC$ 上的射影为 $\triangle ABC$ 的外心 $O$。
	\end{description}

	\textbf{\textcolor{CStatement}{结论}}
	\begin{description}[style=nextline, leftmargin=2.8em, labelsep=0.5em, itemsep=0.45em]
		\item[\textbf{\textcolor{CStatement}{结论一（球心位置）：}}]
		      \textbf{\textcolor{CStatement}{直观描述：}} 外接球球心位于过底面外心 $O$ 且垂直于底面的直线上，特殊情况下可与 $O$ 重合。
		      \[
			      Q \in l,\quad l \perp (ABC),\quad O \in l .
		      \]

		\item[\textbf{\textcolor{CStatement}{结论二（等距条件）：}}]
		      \textbf{\textcolor{CStatement}{直观描述：}} 球心到四个顶点距离都等于外接球半径。
		      \[
			      QA = QB = QC = QS = R .
		      \]

		\item[\textbf{\textcolor{CStatement}{常见变形（半径公式）：}}]
		      \textbf{\textcolor{CStatement}{直观描述：}} 设底面 $\triangle ABC$ 的外接圆半径为 $r$，三棱锥的高为 $h=SO>0$，则外接球半径为
		      \[
			      R = \frac{h^{2} + r^{2}}{2h}\quad (h > 0).
		      \]
	\end{description}

\end{statementbox}